% LEGACY: the elsarticle root used before 2026-10-06, kept for reference only.
% The live root is paper/neurocomputing.tex on the Elsevier CAS class (cas-sc).
% Neurocomputing (Elsevier) build root, 2026-10-06. The SHORT paper: a
% separate body in nc/, written for this submission. It does not input
% main.tex, which stays the long version (the infofus/jair/tmlr roots) and is
% not edited by this build.
%
% Decisions behind it are recorded in notes and memory: the venue (free,
% subscription route), the title (the user's choice), and leaving the
% rate-distortion material out entirely (the user's choice; it lives in the
% long version and the Zenodo preprint).
%
% Neurocomputing asks for: abstract <= 250 words, 3-5 highlights of about 85
% characters, keywords, and a <= 100-word biography plus a passport-type
% photograph for EVERY author, submitted as separate files (nc/author_bios.md).
%
% Build from paper/:  pdflatex neurocomputing && bibtex neurocomputing &&
% pdflatex neurocomputing && pdflatex neurocomputing
% (nc/build.sh does exactly that). Figures: python nc/make_figures.py.
%
% The xcolor option must be passed before the class for the same reason as in
% every other root: a later [table] request option-clashes silently and
% \cellcolor goes undefined.
\PassOptionsToPackage{table}{xcolor}
% Numbered citations in square brackets, ordered by first appearance: the
% Neurocomputing reference style. (At first submission any consistent style is
% accepted under Elsevier's "Your Paper Your Way"; matching the journal's own
% style anyway saves a step at proof.)
\documentclass[preprint,12pt,number,sort&compress]{elsarticle}

\journal{Neurocomputing}

\input{preamble}

% pdfLaTeX only, so microtype's protrusion and expansion are both available
% (the preamble's note against microtype concerns tectonic/XeTeX). xurl lets
% the repository URL, which is long and hyphen-heavy, break at any character.
\usepackage{microtype}
\usepackage{xurl}
\usepackage{tabularx}

% Six full-width multi-panel figures: the default float fractions push them
% all to the end of the document, so allow taller top floats and fuller
% float pages.
\renewcommand{\topfraction}{0.9}
\renewcommand{\bottomfraction}{0.8}
\renewcommand{\textfraction}{0.07}
\renewcommand{\floatpagefraction}{0.75}

\begin{document}

\begin{frontmatter}

\title{When LoRA Adapters Share an Initialisation: Subspace Collapse and Its
Consequences for Merging}

\author[jiet]{Sankalp Pathak\corref{cor1}}
\ead{pathaksankalp04@gmail.com}
\author[jiet]{Sanjay Garg}
\ead{gargsv@gmail.com}
\author[jiet]{Piyush Kumar Singh}
\ead{Piyushsingh5629@gmail.com}
\cortext[cor1]{Corresponding author.}
\affiliation[jiet]{organization={Department of Computer Science and Engineering,
    Jaypee University of Engineering and Technology},
  city={Guna}, state={Madhya Pradesh}, country={India}}

\begin{abstract}
\input{nc/abstract}
\end{abstract}

\begin{highlights}
\input{nc/highlights}
\end{highlights}

\begin{keyword}
LoRA \sep model merging \sep parameter-efficient fine-tuning \sep
initialisation \sep conditioning \sep pre-registration
\end{keyword}

\end{frontmatter}

\section{Introduction}\label{sec:intro}
\input{nc/s1_intro}

\section{Background}\label{sec:background}
\input{nc/s2_background}

\section{Shared initialisation collapses the task subspaces}\label{sec:collapse}
\input{nc/s3_collapse}

\section{A first-order theory of the collapse}\label{sec:theory}
\input{nc/s4_theory}

\section{Where the collapse occurs}\label{sec:prevalence}
\input{nc/s5_prevalence}

\section{Which merges the collapse affects}\label{sec:merging}
\input{nc/s6_merging}

\section{Discussion and limitations}\label{sec:discussion}
\input{nc/s7_discussion}

\section{Conclusion}\label{sec:conclusion}
\input{nc/s8_conclusion}

\input{nc/declarations}

\bibliographystyle{elsarticle-num}
\bibliography{references}

\appendix
\section{Derivation of the first-order theory}\label{app:theory}
\input{nc/app_theory}

\section{Pre-registration and commit trail}\label{app:prereg}
\input{nc/app_prereg}

\end{document}
